\scanpage{27}
C'est une $j$-stratégie, \deleted{prouvons} \inserted{Prouvons} qu'elle est gagnante pour $x_0$. \deleted{\unsure{ou}} Une partie maximale jouée suivant $\Sigma_0$ \unsure{commence} avec $(x_0,x_1)$, soit $(x_0,x_1,\ldots,x_m,\ldots)$ \deleted{les $x_i$ $(i\geq1)$ sont \unsure{gagnants} pour $\Sigma_1$}, prouvons qu'elle est gagnée par $j$. Prouvons que tous les $x_i$ $(i\geq1)$ sont $\ne x_0$. Sinon, \deleted{si $i$ est le} plus petit indice $i\geq1$ tel que $x_i=x_0$, \deleted{pour un $\unsure{i>1|j>1}$}, alors \deleted{par la prop. \unsure{il s'ensuivrait}} \inserted{$(x_1,\ldots,x_i=x_0)$ serait compatible avec $(\Sigma_1)$,} on aurait que $x_i\inserted{{}=x_0}$ est gagnant pour $\Sigma_1$, contrairement à la supposition que $x_0$ non gagnant pour $\Sigma_1$.

Comme les $x_i$ $(i\geq1)$ sont $\ne x_0$, la partie $(x_1,x_2,\deleted{\unsure{x_0}}\inserted{\ldots})$ est une $\Sigma_1$-partie, donc elle est gagnante pour $j$, donc $(x_0,x_1,\ldots)$ aussi; cqfd.

b) On suppose $x_0\in\bar C(j)$.

1) $\subset$. Supposons $x_0\in G(j)$, prouvons $R(x)\subset G(j)$, Soit $x_1\in R(x)$, prouvons $x_1\in G(j)$. En effet, \inserted{si $\Sigma_0$ est une $j$-stratégie $x_0$-gagn.}, $(x_0,x_1)$ est une \inserted{$\Sigma_0$-}partie, donc par la prop., $x_1$ est gagnant pour $\Sigma_0$, donc $x_1\in G(j)$.

2) $\supset$. Supposons $R(x_0)\subset G(j)$, prouvons $x_0\in G(j)$. $\forall x_1\in R(x_0)$, soit $\Sigma_{x_1}$ une $j$-stratégie $x_1$-gagnante. Choisissons un \uline{bon ordre} sur $R(x_0)$. On va définir une $j$-stratégie $\Sigma_0$ $x_0$-gagnante ainsi :

$\forall x\in C$, soit
\[
E(x)=\{x_1\in R(x_0)\mid\exists\ \Sigma_{x_1}\text{-partie de }\unsure{x_0|x_1}\text{ à }x\}
\]
% archive foliation: 26
